
\documentclass{article}
\usepackage{colortbl}
\usepackage{makecell}
\usepackage{multirow}
\usepackage{supertabular}

\begin{document}

\newcounter{utterance}

\centering \large Interaction Transcript for game `eqbench', experiment `eq\_bench', episode 54 with reflexion{-}Qwen3.5{-}27B.
\vspace{24pt}

{ \footnotesize  \setcounter{utterance}{1}
\setlength{\tabcolsep}{0pt}
\begin{supertabular}{c@{$\;$}|p{.15\linewidth}@{}p{.15\linewidth}p{.15\linewidth}p{.15\linewidth}p{.15\linewidth}p{.15\linewidth}}
    \# & \multicolumn{2}{c}{Player} && \multicolumn{2}{c}{Game Master} \\
    \hline

    \theutterance \stepcounter{utterance}  
    & & & \multicolumn{4}{p{0.6\linewidth}}{
        \cellcolor[rgb]{0.9,0.9,0.9}{
            \makecell[{{p{\linewidth}}}]{
                \texttt{\tiny{[P1$\langle$GM]}}
                \texttt{You have attempted this same task before and did not succeed. These are the plans you wrote after those attempts.} \\
\texttt{Trial \#1: New plan: My previous attempt failed because I assigned a non{-}zero score to "Insecure" (4), which was likely incorrect given the context. Reggie's response ("What's that supposed to mean?") and his confident retort ("I can be full of surprises") indicate he is reacting defensively to sarcasm rather than feeling self{-}doubt. While "Annoyed" is the primary emotion, "Insecure" should be scored as 0 because Reggie maintains a defensive, confident posture rather than a vulnerable one. Next time, I will strictly evaluate whether the character's dialogue supports a specific emotion before assigning a score > 0, ensuring that emotions like "Insecure" are only scored if there is clear evidence of self{-}doubt in the text. I will also double{-}check that the sum of scores reflects the intensity hierarchy correctly, likely keeping "Annoyed" high and "Insecure" at 0.} \\
\\ 
\texttt{Trial \#2: New plan: My previous attempt failed because I assigned a score of 0 to "Insecure," missing the subtle emotional undercurrent in the dialogue. While Reggie's outward reaction is defensive and annoyed ("What's that supposed to mean?"), Marcus's sarcasm about his "powerpoint presentation" and "five{-}year plan" directly attacks his competence and social style, which typically triggers insecurity in such contexts. By scoring "Insecure" as 0, I likely failed to capture the full emotional complexity required for the target state. Next time, I will analyze not just the explicit defensive tone but also the implicit vulnerability caused by the other character's mockery. I will assign a non{-}zero score (e.g., 4 or 5) to "Insecure" to reflect the sting of the personal jab, while keeping "Annoyed" as the highest score (e.g., 8 or 9) to represent the immediate reaction. I will ensure that at least two emotions are scored if the context supports a mix of surface anger and underlying doubt.} \\
\\ 
\\ 
\\ 
\texttt{Your task is to predict the likely emotional responses of a character in this dialogue:} \\
\\ 
\texttt{Reggie: Marcus, I've been thinking about Mom's birthday party. I've got some ideas.} \\
\texttt{Marcus: Oh, I bet you do. Another corporate gala, Reggie? } \\
\texttt{Reggie: Funny. I was actually thinking of something more... intimate. } \\
\texttt{Marcus: Intimate? You? Now, that's a surprise.} \\
\texttt{Reggie: Yes, well, I can be full of surprises. } \\
\texttt{Marcus: I'll believe it when I see it. } \\
\texttt{Reggie: What's that supposed to mean? } \\
\texttt{Marcus: Nothing. Just that your idea of surprise usually involves a powerpoint presentation and a five{-}year plan.} \\
\texttt{{[}End dialogue{]}} \\
\\ 
\texttt{At the end of this dialogue, Reggie would feel...} \\
\texttt{Vindicated} \\
\texttt{Annoyed} \\
\texttt{Insecure} \\
\texttt{Euphoric} \\
\\ 
\texttt{Give each of these possible emotions a score from 0{-}10 for the relative intensity that they are likely to be feeling each.} \\
\\ 
\texttt{You must output in the following format, including headings (of course, you should give your own scores), with no additional commentary:} \\
\\ 
\texttt{Vindicated: <score>} \\
\texttt{Annoyed: <score>} \\
\texttt{Insecure: <score>} \\
\texttt{Euphoric: <score>} \\
\\ 
\\ 
\texttt{{[}End of answer{]}} \\
\\ 
\texttt{Remember: zero is a valid score, meaning they are likely not feeling that emotion. You must score at least one emotion > 0.} \\
\\ 
\texttt{Your answer:} \\
            }
        }
    }
     \\ \\

    \theutterance \stepcounter{utterance}  
    & \multicolumn{4}{p{0.6\linewidth}}{
        \cellcolor[rgb]{0.561,0.737,0.561}{
            \makecell[{{p{\linewidth}}}]{
                \texttt{\tiny{[P1$\rangle$GM]}}
                \texttt{Vindicated: 0} \\
\texttt{Annoyed: 9} \\
\texttt{Insecure: 4} \\
\texttt{Euphoric: 0} \\
            }
        }
    }
    & & \\ \\

    \theutterance \stepcounter{utterance}  
    & & & \multicolumn{2}{p{0.3\linewidth}}{
        \cellcolor[rgb]{0.9,0.9,0.9}{
            \makecell[{{p{\linewidth}}}]{
                \texttt{\tiny{[GM$|$GM]}}
                \texttt{Vindicated: 0} \\
\texttt{Annoyed: 9} \\
\texttt{Insecure: 4} \\
\texttt{Euphoric: 0} \\
            }
        }
    }
    & & \\ \\

    \theutterance \stepcounter{utterance}  
    & & & \multicolumn{2}{p{0.3\linewidth}}{
        \cellcolor[rgb]{0.9,0.9,0.9}{
            \makecell[{{p{\linewidth}}}]{
                \texttt{\tiny{[GM$|$GM]}}
                \texttt{Vindicated: 0} \\
\texttt{Annoyed: 7} \\
\texttt{Insecure: 4} \\
\texttt{Euphoric: 0} \\
            }
        }
    }
    & & \\ \\

    \theutterance \stepcounter{utterance}  
    & & & \multicolumn{2}{p{0.3\linewidth}}{
        \cellcolor[rgb]{0.9,0.9,0.9}{
            \makecell[{{p{\linewidth}}}]{
                \texttt{\tiny{[GM$|$GM]}}
                \texttt{game\_result = LOSE} \\
            }
        }
    }
    & & \\ \\

    \theutterance \stepcounter{utterance}  
    & & \multicolumn{4}{p{0.6\linewidth}}{
        \cellcolor[rgb]{0.95,0.95,0.95}{
            \makecell[{{p{\linewidth}}}]{
                \texttt{\tiny{[P1$\leftrightarrow$P]}}
                \texttt{Post{-}episode Reflexion (not sent to GM):} \\
\texttt{New plan: My previous attempts failed because I oscillated between two contradictory interpretations of Reggie's emotional state: first assuming he was purely defensive (scoring Insecure 0), and then assuming the sarcasm necessarily implied deep vulnerability (scoring Insecure 4). The dialogue shows Reggie reacting with a sharp, defensive retort ("I can be full of surprises") immediately after being mocked. This suggests his primary reaction is irritation at the interruption and the tone, rather than a collapse into self{-}doubt. The "Insecure" score should likely be 0 or very low because he does not retreat or apologize; he counters. The "Annoyed" score should be the dominant factor. Next time, I will prioritize the character's *active response* over the *intent of the insult*. If the character fights back confidently, I will score "Insecure" as 0. I will assign a high score (8{-}9) to "Annoyed" and 0 to "Vindicated" and "Euphoric" to reflect the clear negative tone without over{-}interpreting hidden vulnerability.} \\
            }
        }
    }
    & \\ \\

\end{supertabular}
}

\end{document}
